\begin{proof}[Proof of \cref{obj:lem:fixed-cube-holder-completion}]
\textbf{1. Parameters and derivative notation.}
Set
\[
s:=\beta-m,\qquad
\mathscr C:=[-1,1]^d,\qquad
\mathscr C^\circ:=(-1,1)^d,\qquad
\mathscr U:=[0,1]^d.
\]
Since \(m=\lceil\beta\rceil-1\) and \(\beta>1\),
\[
m\ge1,\qquad
m<\beta\le m+1,\qquad
0<s\le1.
\]
For each integer \(j\ge0\), define the set of ordered coordinate indices and the associated derivative by
\[
\mathscr I_j:=\{1,\ldots,d\}^{\{1,\ldots,j\}},
\qquad
\partial_\iota^j w
:=
\partial_{\iota(1)}\cdots\partial_{\iota(j)}w,
\quad \iota\in\mathscr I_j.
\]
Here \(\mathscr I_0\) contains the empty index and
\(\partial_\iota^0w=w\). On a closed cube, these derivatives use their continuous interior traces. Symmetry of the interior derivatives identifies an ordered derivative with the multi-index derivative having the same coordinate multiplicities; continuity preserves this identity on the boundary.
% lean: CausalSmith.Stat.WeakOverlap.fixedCube_holder_completion

\textbf{2. Completion on the source cube.}
We first derive a constant \(K_0>0\), depending only on \(d,m,s\), with the following property. For arbitrary
\[
u_{\mathrm{aux}}:\mathbb R^d\to\mathbb R,
\qquad
M_{\mathrm{aux}},L_{\mathrm{aux}}\in[0,\infty),
\]
suppose that
\[
u_{\mathrm{aux}}\in C^m(\mathscr C),\qquad
|u_{\mathrm{aux}}(x)|\le M_{\mathrm{aux}}
\quad(x\in\mathscr C),
\]
and
\[
|\partial_\iota^m u_{\mathrm{aux}}(x)
-\partial_\iota^m u_{\mathrm{aux}}(y)|
\le L_{\mathrm{aux}}\|x-y\|_\infty^s
\quad(\iota\in\mathscr I_m,\ x,y\in\mathscr C).
\]
The desired completion estimates are
\[
|\partial_\iota^j u_{\mathrm{aux}}(x)|
\le K_0(M_{\mathrm{aux}}+L_{\mathrm{aux}})
\quad(0\le j\le m,\ \iota\in\mathscr I_j,\ x\in\mathscr C)
\]
and
\[
|\partial_\iota^m u_{\mathrm{aux}}(x)
-\partial_\iota^m u_{\mathrm{aux}}(y)|
\le K_0(M_{\mathrm{aux}}+L_{\mathrm{aux}})
\|x-y\|_\infty^s.
\]

To obtain the interpolation estimate used in this completion, first consider a function whose ordinary coordinate derivatives on \(\mathscr C\) satisfy the displayed top-order modulus. Multilinearity gives, for a direction vector \(v\in\mathbb R^d\),
\[
D^m u_{\mathrm{aux}}(x)[v,\ldots,v]
=
\sum_{\iota\in\mathscr I_m}
\left(\prod_{a=1}^m v_{\iota(a)}\right)
\partial_\iota^m u_{\mathrm{aux}}(x).
\]
Consequently,
\[
\begin{aligned}
&|D^m u_{\mathrm{aux}}(x)[v,\ldots,v]
-D^m u_{\mathrm{aux}}(y)[v,\ldots,v]|\\
&\qquad\le
d^m L_{\mathrm{aux}}\|x-y\|_\infty^s\|v\|_\infty^m\\
&\qquad\le
(d+1)^m L_{\mathrm{aux}}\|x-y\|_\infty^s\|v\|_\infty^m.
\end{aligned}
\]
Define
\[
C_{\mathrm{diag}}:=(d+1)^m,\qquad
\varphi_{x,y}(r):=
u_{\mathrm{aux}}\bigl(x+r(y-x)\bigr),
\quad r\in\mathbb R,
\]
for \(x\in\mathscr C^\circ\) and \(y\in\mathscr C\).
The segment lies in \(\mathscr C\), and its points with \(0\le r<1\) lie in \(\mathscr C^\circ\). The chain rule therefore gives
\[
\varphi_{x,y}^{(k)}(r)
=
D^k u_{\mathrm{aux}}\bigl(x+r(y-x)\bigr)
[y-x,\ldots,y-x],
\qquad 0\le k\le m,\quad 0\le r<1.
\]
Using the diagonal estimate and
\[
\|x+r(y-x)-x-r'(y-x)\|_\infty
=
|r-r'|\|y-x\|_\infty,
\]
we obtain
\[
|\varphi_{x,y}^{(m)}(r)-\varphi_{x,y}^{(m)}(r')|
\le
C_{\mathrm{diag}}L_{\mathrm{aux}}
\|y-x\|_\infty^{m+s}|r-r'|^s,
\qquad r,r'\in[0,1).
\]
Taylor's Lagrange remainder of degree \(m-1\), applied on \([0,1]\), supplies \(\xi\in(0,1)\) such that
\[
\varphi_{x,y}(1)
-\sum_{k=0}^{m-1}\frac{\varphi_{x,y}^{(k)}(0)}{k!}
=
\frac{\varphi_{x,y}^{(m)}(\xi)}{m!}.
\]
After subtracting the degree-\(m\) term,
\[
\varphi_{x,y}(1)
-\sum_{k=0}^{m}\frac{\varphi_{x,y}^{(k)}(0)}{k!}
=
\frac{\varphi_{x,y}^{(m)}(\xi)
-\varphi_{x,y}^{(m)}(0)}{m!}.
\]
Since \(\xi^s\le1\) and \(m!\ge1\), this proves
\[
\left|
u_{\mathrm{aux}}(y)
-\sum_{k=0}^{m}\frac{1}{k!}
D^k u_{\mathrm{aux}}(x)[y-x,\ldots,y-x]
\right|
\le
C_{\mathrm{diag}}L_{\mathrm{aux}}
\|y-x\|_\infty^{m+s}.
\]

We next bound this Taylor polynomial on a fixed tensor grid. Define
\[
\nu_j:=\frac{j}{m+1}\quad(0\le j\le m),\qquad
\mathscr A:=\{0,\ldots,m\}^d,\qquad
z_b:=(\nu_{b_1},\ldots,\nu_{b_d})
\quad(b\in\mathscr A).
\]
Every \(z_b\) belongs to \(\mathscr C\). Compactness supplies a finite number \(D_{\mathrm{cube}}\ge0\) satisfying
\[
\|y-x\|_\infty\le D_{\mathrm{cube}}
\quad(x,y\in\mathscr C).
\]
Set
\[
\kappa_{\mathrm{cube}}:=\max\{1,D_{\mathrm{cube}}\},
\qquad
C_{\mathrm{grid}}
:=
1+C_{\mathrm{diag}}\kappa_{\mathrm{cube}}^{m+s}.
\]
For the Taylor polynomial
\[
p_x(z):=
\sum_{k=0}^{m}\frac{1}{k!}
D^k u_{\mathrm{aux}}(x)[z-x,\ldots,z-x],
\qquad z\in\mathbb R^d,
\]
the response bound and remainder estimate yield
\[
\begin{aligned}
|p_x(z_b)|
&\le |u_{\mathrm{aux}}(z_b)|
   +|u_{\mathrm{aux}}(z_b)-p_x(z_b)|\\
&\le M_{\mathrm{aux}}
   +C_{\mathrm{diag}}L_{\mathrm{aux}}
      \kappa_{\mathrm{cube}}^{m+s}\\
&\le C_{\mathrm{grid}}
      (M_{\mathrm{aux}}+L_{\mathrm{aux}}).
\end{aligned}
\]

Here is the polynomial estimate that converts these grid bounds into derivative bounds. Define the univariate Lagrange polynomials, their coefficients, and the tensor weights by
\[
\Lambda_j(t):=
\prod_{\substack{0\le q\le m\\q\ne j}}
\frac{t-\nu_q}{\nu_j-\nu_q},
\qquad
\omega_{kj}:=[t^k]\Lambda_j(t),
\qquad
W_{ab}:=\prod_{i=1}^d\omega_{a_i b_i},
\]
where \(0\le k,j\le m\) and \(a,b\in\mathscr A\).
For any polynomial \(p_{\mathrm{aux}}\) of degree at most \(m\) in each coordinate, define
\[
c_a(p_{\mathrm{aux}}):=[z^a]p_{\mathrm{aux}}(z),
\qquad
z^a:=\prod_{i=1}^d z_i^{a_i}.
\]
Univariate Lagrange interpolation, applied to each monomial, gives
\[
\sum_{j=0}^m\omega_{kj}\nu_j^q
=
\begin{cases}
1,&k=q,\\
0,&k\ne q,
\end{cases}
\qquad 0\le k,q\le m.
\]
Taking products over the coordinates and then expanding \(p_{\mathrm{aux}}\) proves the coefficient recovery identity
\[
c_a(p_{\mathrm{aux}})
=
\sum_{b\in\mathscr A}W_{ab}p_{\mathrm{aux}}(z_b).
\]

The weights admit a uniform bound. Define the coefficient one-norm of a univariate polynomial by
\[
\|p_{\mathrm{uni}}\|_{\mathrm{coef},1}
:=
\sum_{k\ge0}|[t^k]p_{\mathrm{uni}}(t)|.
\]
The coefficient convolution formula and the triangle inequality imply
\[
\|p_{\mathrm{uni}}q_{\mathrm{uni}}\|_{\mathrm{coef},1}
\le
\|p_{\mathrm{uni}}\|_{\mathrm{coef},1}
\|q_{\mathrm{uni}}\|_{\mathrm{coef},1}.
\]
Since \(0\le\nu_q\le1\), each numerator factor has coefficient one-norm at most \(2\). Splitting the denominator factors below and above \(j\) gives
\[
\left|
\prod_{\substack{0\le q\le m\\q\ne j}}(\nu_j-\nu_q)
\right|
=
\frac{j!(m-j)!}{(m+1)^m}.
\]
Thus
\[
\sum_{k=0}^m|\omega_{kj}|
=
\|\Lambda_j\|_{\mathrm{coef},1}
\le
\frac{2^m(m+1)^m}{j!(m-j)!}.
\]
The binomial identity and the elementary binomial coefficient bounds give
\[
\sum_{j=0}^m\frac{1}{j!(m-j)!}
=\frac{2^m}{m!},
\qquad
\frac{(m+1)^m}{m!}
\le\binom{2m}{m}\le 2^{2m}=4^m.
\]
Consequently,
\[
\sum_{k=0}^m\sum_{j=0}^m|\omega_{kj}|
\le
4^m\frac{(m+1)^m}{m!}
\le 2^{4m}
\le 2^{4m+4}.
\]
Define
\[
C_{\mathrm{coef}}:=(2^{4m+4})^d.
\]
Factoring the finite tensor sums now yields
\[
\sum_{a\in\mathscr A}\sum_{b\in\mathscr A}|W_{ab}|
=
\left(\sum_{k=0}^m\sum_{j=0}^m|\omega_{kj}|\right)^d
\le C_{\mathrm{coef}}.
\]
Hence, whenever \(B_{\mathrm{pol}}\ge0\) and
\(|p_{\mathrm{aux}}(z_b)|\le B_{\mathrm{pol}}\) for every \(b\in\mathscr A\),
\[
|c_a(p_{\mathrm{aux}})|
\le
\sum_{b\in\mathscr A}|W_{ab}|
|p_{\mathrm{aux}}(z_b)|
\le C_{\mathrm{coef}}B_{\mathrm{pol}}.
\]

For each \(a\in\mathscr A\), the monomial \(z\mapsto z^a\) is smooth. Continuity of its iterated differentials on the compact cube, followed by a maximum over the finitely many \(a\) and \(0\le j\le m\), supplies \(M_{\mathrm{mon}}\ge0\) such that
\[
\|D^j(z\mapsto z^a)(x)\|_{\mathrm{op},\infty}
\le M_{\mathrm{mon}}
\quad(a\in\mathscr A,\ 0\le j\le m,\ x\in\mathscr C).
\]
Each coordinate unit vector has supremum norm \(1\), so
\[
|\partial_\iota^j(z\mapsto z^a)(x)|
\le M_{\mathrm{mon}}.
\]
Define
\[
C_{\mathrm{pol}}
:=
(m+1)^d C_{\mathrm{coef}}(M_{\mathrm{mon}}+1).
\]
Differentiating the finite expansion
\(p_{\mathrm{aux}}(z)=\sum_{a\in\mathscr A}c_a(p_{\mathrm{aux}})z^a\)
therefore gives
\[
\begin{aligned}
|\partial_\iota^j p_{\mathrm{aux}}(x)|
&\le
\sum_{a\in\mathscr A}|c_a(p_{\mathrm{aux}})|
|\partial_\iota^j(z\mapsto z^a)(x)|\\
&\le (m+1)^d C_{\mathrm{coef}}B_{\mathrm{pol}}M_{\mathrm{mon}}\\
&\le C_{\mathrm{pol}}B_{\mathrm{pol}},
\end{aligned}
\]
for \(0\le j\le m\), \(\iota\in\mathscr I_j\), and \(x\in\mathscr C\).

The polynomial \(p_x\) has total degree at most \(m\), as is also seen from its coordinate expansion
\[
p_x(z)=
\sum_{k=0}^m\frac{1}{k!}
\sum_{\iota\in\mathscr I_k}
\partial_\iota^k u_{\mathrm{aux}}(x)
\prod_{a=1}^k(z_{\iota(a)}-x_{\iota(a)}).
\]
Its derivatives at the expansion point agree with those of \(u_{\mathrm{aux}}\). Indeed, a homogeneous term of degree \(k<j\) has zero \(j\)-th derivative; a term of degree \(k>j\) retains a vanishing factor at \(x\); and for \(k=j\), symmetry of the differential produces \(j!\) identical contributions, cancelling \(1/j!\). Thus
\[
\partial_\iota^j p_x(x)
=
\partial_\iota^j u_{\mathrm{aux}}(x),
\qquad 0\le j\le m.
\]
Set
\[
K_{\mathrm{int}}:=C_{\mathrm{pol}}C_{\mathrm{grid}}>0.
\]
Combining the grid bound, polynomial derivative bound, and derivative agreement proves
\[
|\partial_\iota^j u_{\mathrm{aux}}(x)|
\le K_{\mathrm{int}}(M_{\mathrm{aux}}+L_{\mathrm{aux}})
\quad(x\in\mathscr C^\circ)
\]
under the ordinary-derivative modulus used above.

To apply this estimate to intrinsic cube derivatives, define the inward contractions
\[
u_{\mathrm{aux},t}(z):=u_{\mathrm{aux}}(tz),
\qquad z\in\mathbb R^d,\quad 0<t<1.
\]
For every \(z\in\mathscr C\), \(tz\in\mathscr C^\circ\). The chain rule on this interior region gives
\[
u_{\mathrm{aux},t}\in C^m(\mathscr C),\qquad
|u_{\mathrm{aux},t}(z)|\le M_{\mathrm{aux}},
\qquad
\partial_\iota^j u_{\mathrm{aux},t}(z)
=t^j\partial_\iota^j u_{\mathrm{aux}}(tz).
\]
The assumed intrinsic top-order modulus implies
\[
\begin{aligned}
&|\partial_\iota^m u_{\mathrm{aux},t}(x)
-\partial_\iota^m u_{\mathrm{aux},t}(y)|\\
&\qquad\le
t^mL_{\mathrm{aux}}\|tx-ty\|_\infty^s\\
&\qquad=
t^m t^s L_{\mathrm{aux}}\|x-y\|_\infty^s
\le L_{\mathrm{aux}}\|x-y\|_\infty^s,
\end{aligned}
\]
because \(t^m\le1\) and \(t^s\le1\). The preceding interpolation estimate therefore applies to \(u_{\mathrm{aux},t}\), yielding
\[
|t^j\partial_\iota^j u_{\mathrm{aux}}(tx)|
\le K_{\mathrm{int}}(M_{\mathrm{aux}}+L_{\mathrm{aux}})
\quad(x\in\mathscr C^\circ,\ 0<t<1).
\]
Continuity of the derivative traces allows \(t\uparrow1\), giving
\[
|\partial_\iota^j u_{\mathrm{aux}}(x)|
\le K_{\mathrm{int}}(M_{\mathrm{aux}}+L_{\mathrm{aux}})
\quad(x\in\mathscr C^\circ).
\]
Since \(\overline{\mathscr C^\circ}=\mathscr C\), the same continuity extends this bound to every \(x\in\mathscr C\).

Finally, define
\[
K_0:=K_{\mathrm{int}}+1>0.
\]
Nonnegativity of the bounds gives
\[
K_{\mathrm{int}}(M_{\mathrm{aux}}+L_{\mathrm{aux}})
\le K_0(M_{\mathrm{aux}}+L_{\mathrm{aux}}),
\qquad
L_{\mathrm{aux}}
\le K_0(M_{\mathrm{aux}}+L_{\mathrm{aux}}).
\]
These inequalities prove both completion estimates stated at the beginning of this step, together with the original \(C^m\) regularity.
% lean: Causalean.Mathlib.Analysis.Calculus.CubeExtension.cube_holder_completion

\textbf{3. Affine transport of a completed radius.}
Define
\[
K_1:=2^{m+1}>0.
\]
Suppose that a function \(u_{\mathrm{aux}}\in C^m(\mathscr C)\) has derivative and top-order modulus bounds of radius \(\rho\ge0\):
\[
|\partial_\iota^j u_{\mathrm{aux}}(x)|\le\rho,
\qquad
|\partial_\iota^m u_{\mathrm{aux}}(x)
-\partial_\iota^m u_{\mathrm{aux}}(y)|
\le\rho\|x-y\|_\infty^s.
\]
The affine map \(T(x)=2x-\mathbf1\) maps \(\mathscr U\) into \(\mathscr C\). Composition preserves the intrinsic \(C^m\) regularity, and the chain rule gives
\[
\partial_\iota^j(u_{\mathrm{aux}}\circ T)(x)
=
2^j\partial_\iota^j u_{\mathrm{aux}}(T(x)),
\qquad x\in\mathscr U,\quad 0\le j\le m.
\]
This identity holds on the interior and extends to the boundary by the continuous traces. Hence
\[
|\partial_\iota^j(u_{\mathrm{aux}}\circ T)(x)|
\le2^j\rho\le K_1\rho.
\]
Moreover,
\[
\|T(x)-T(y)\|_\infty=2\|x-y\|_\infty,
\qquad 2^s\le2,
\]
so
\[
\begin{aligned}
&|\partial_\iota^m(u_{\mathrm{aux}}\circ T)(x)
-\partial_\iota^m(u_{\mathrm{aux}}\circ T)(y)|\\
&\qquad\le
2^m\rho\|T(x)-T(y)\|_\infty^s\\
&\qquad=
2^m2^s\rho\|x-y\|_\infty^s\\
&\qquad\le
K_1\rho\|x-y\|_\infty^s.
\end{aligned}
\]
Thus affine transport multiplies the completed radius by at most \(K_1\).
% lean: Causalean.Mathlib.Analysis.Calculus.CubeExtension.cube_holder_affine_transport

\textbf{4. Conversion of the given modulus.}
Choose
\[
C_{\mathrm{norm}}:=(\sqrt d)^s>0,
\qquad
K_{d,\beta}:=K_1K_0(C_{\mathrm{norm}}+1)>0.
\]
These constants depend only on \(d\) and \(\beta\). For any
\(\iota\in\mathscr I_m\), define its coordinate multiplicities by
\[
\alpha_i:=\#\{a\in\{1,\ldots,m\}:\iota(a)=i\},
\qquad i=1,\ldots,d.
\]
Then
\[
|\alpha|=\sum_{i=1}^d\alpha_i=m,
\qquad
\partial_\iota^m u=D^\alpha u.
\]
The hypothesis therefore bounds every ordered top derivative. For \(x,y\in\mathscr C\),
\[
\sum_{i=1}^d(x_i-y_i)^2
\le d\|x-y\|_\infty^2,
\qquad
\left(\sum_{i=1}^d(x_i-y_i)^2\right)^{1/2}
\le\sqrt d\,\|x-y\|_\infty.
\]
Since \(s>0\), raising the latter inequality to the power \(s\) gives
\[
\left(\sum_{i=1}^d(x_i-y_i)^2\right)^{s/2}
\le C_{\mathrm{norm}}\|x-y\|_\infty^s.
\]
Multiplication by \(L_0\ge0\) now yields
\[
|\partial_\iota^m u(x)-\partial_\iota^m u(y)|
\le
C_{\mathrm{norm}}L_0\|x-y\|_\infty^s.
\]
% lean: CausalSmith.Stat.WeakOverlap.fixedCube_holder_completion

\textbf{5. Completion and transport for \(u\).}
Define the source radius
\[
\rho_0:=K_0(M_0+C_{\mathrm{norm}}L_0)\ge0.
\]
The source-cube completion estimates apply with the response bound \(M_0\) and modulus constant \(C_{\mathrm{norm}}L_0\). They give
\[
u\in C^m(\mathscr C),\qquad
|\partial_\iota^j u(x)|\le\rho_0
\quad(0\le j\le m,\ \iota\in\mathscr I_j,\ x\in\mathscr C),
\]
and
\[
|\partial_\iota^m u(x)-\partial_\iota^m u(y)|
\le\rho_0\|x-y\|_\infty^s.
\]
Applying the affine transport estimate with radius \(\rho_0\) proves
\[
u\circ T\in C^m(\mathscr U),\qquad
|\partial_\iota^j(u\circ T)(x)|\le K_1\rho_0,
\]
and
\[
|\partial_\iota^m(u\circ T)(x)
-\partial_\iota^m(u\circ T)(y)|
\le K_1\rho_0\|x-y\|_\infty^s.
\]
% lean: CausalSmith.Stat.WeakOverlap.fixedCube_holder_completion

\textbf{6. Calibration of the final radius.}
The nonnegative bounds imply
\[
\begin{aligned}
&(C_{\mathrm{norm}}+1)(M_0+L_0)
-(M_0+C_{\mathrm{norm}}L_0)\\
&\qquad=C_{\mathrm{norm}}M_0+L_0\ge0.
\end{aligned}
\]
Multiplying by \(K_1K_0>0\) gives
\[
K_1\rho_0
=
K_1K_0(M_0+C_{\mathrm{norm}}L_0)
\le K_{d,\beta}(M_0+L_0).
\]
Consequently,
\[
\max_{|\alpha|\le m}
\|D^\alpha(u\circ T)\|_\infty
\le K_{d,\beta}(M_0+L_0),
\]
and, for every \(|\alpha|=m\) and \(x,y\in\mathscr U\),
\[
|D^\alpha(u\circ T)(x)-D^\alpha(u\circ T)(y)|
\le
K_{d,\beta}(M_0+L_0)\|x-y\|_\infty^{\beta-m}.
\]
Together with the transported intrinsic \(C^m\) regularity, these are precisely the restrictions in
\cref{obj:ass:holder-derivative-bound,obj:ass:holder-modulus}.
By \cref{obj:def:holder},
\[
u\circ T\in
\mathcal H^\beta\!\left(K_{d,\beta}(M_0+L_0)\right).
\]
All estimates include zero values of either bound, and the argument uses \(0<s\le1\), including \(s=1\).
% lean: CausalSmith.Stat.WeakOverlap.fixedCube_holder_completion
\end{proof}
